\documentclass[../../main.tex]{subfiles}

\begin{document}

\section{和、积与复合}\label{section:ch03-和积复合}

本章考虑的所有函数取值于扩充实数, 即集合 $\mathbf{R}\cup\{\pm\infty\}$. 回忆一下, 一个性质称为在可测集 $E$ 上几乎处处（简写为 a.e.）成立, 若它在 $E\sim E_{0}$ 上成立, 其中 $E_{0}$ 是 $E$ 的满足 $m(E_{0})=0$ 的子集.

给定定义在 $E$ 上的两个函数 $h$ 和 $g$. 为了记号上的简洁, 我们常写 "在 $E$ 上, $h\leqslant g$" 来表示对所有 $x\in E$, $h(x)\leqslant g(x)$. 我们说 $E$ 上的函数序列 $\{f_{n}\}$ 是递增的, 若对每个指标 $n$, 在 $E$ 上 $f_{n}\leqslant f_{n+1}$.

\begin{proposition}\label{proposition:ch03-可测等价}
令函数 $f$ 有可测定义域 $E$. 则以下叙述等价:
\begin{enumerate}[(i)]
\item 对每个实数 $c$, 集合 $\{x\in E\mid f(x)>c\}$ 是可测的.
\item 对每个实数 $c$, 集合 $\{x\in E\mid f(x)\geqslant c\}$ 是可测的.
\item 对每个实数 $c$, 集合 $\{x\in E\mid f(x)<c\}$ 是可测的.
\item 对每个实数 $c$, 集合 $\{x\in E\mid f(x)\leqslant c\}$ 是可测的.
\end{enumerate}
这些性质中的每一个都蕴涵着对每个扩充实数 $c$, 集合 $\{x\in E\mid f(x)=c\}$ 是可测的.
\end{proposition}

\begin{proof}
如同 (ii) 和 (iii) 中的集合, 由于 (i) 和 (iv) 中的集合在 $E$ 中互为补集, 而 $E$ 的可测子集在 $E$ 中的补集是可测的, 所以 (ii) 和 (iii)、(i) 和 (iv) 是等价的.

现在 (i) 蕴涵 (ii), 这是由于
\begin{align*}
\{x\in E\mid f(x)\geqslant c\}=\bigcap_{k=1}^{\infty}\{x\in E\mid f(x)>c-1/k\},
\end{align*}
可测集的可数族的交是可测的. 类似地, (ii) 蕴涵 (i), 这是由于
\begin{align*}
\{x\in E\mid f(x)>c\}=\bigcup_{k=1}^{\infty}\{x\in E\mid f(x)\geqslant c+1/k\},
\end{align*}
可测集的可数族的并是可测的.

因此陈述 (i)$\sim$(iv) 是等价的. 现在假设它们中的一个成立, 因此它们中的所有成立. 若 $c$ 是实数, $\{x\in E\mid f(x)=c\}=\{x\in E\mid f(x)\geqslant c\}\cap\{x\in E\mid f(x)\leqslant c\}$, 因此 $f^{-1}(c)$ 是可测的, 这是由于它是两个可测集的交. 另一方面, 若 $c$ 是无穷的, 即 $c=\infty$,
\begin{align*}
\{x\in E\mid f(x)=\infty\}=\bigcap_{k=1}^{\infty}\{x\in E\mid f(x)>k\},
\end{align*}
因此 $f^{-1}(\infty)$ 是可测的, 这是由于它是可测集的可数族的交.
\end{proof}

\begin{definition}[Lebesgue 可测函数]\label{definition:ch03-可测函数}
定义在 $E$ 上的扩充的实值函数 $f$ 称为是 Lebesgue 可测的, 或简单地称为可测的, 若它的定义域 $E$ 是可测的且它满足\cref{proposition:ch03-可测等价}的四个陈述之一.
\end{definition}

\begin{proposition}\label{proposition:ch03-开集原象}
令 $f$ 为定义在可测集 $E$ 上的实值函数. 则函数 $f$ 是可测的当且仅当对每个开集 $\mathcal{O}$, $\mathcal{O}$ 在 $f$ 下的原象 $f^{-1}(\mathcal{O})=\{x\in E\mid f(x)\in\mathcal{O}\}$ 是可测的.
\end{proposition}

\begin{proof}
若每个开集的原象是可测的, 则由于每个区间 $(c,\infty)$ 是开的, 函数 $f$ 是可测的. 反过来, 假定 $f$ 是可测的. 令 $\mathcal{O}$ 为开的. 则我们可将 $\mathcal{O}$ 表示为开有界区间的可数族 $\{I_{k}\}_{k=1}^{\infty}$ 的并, 其中每个 $I_{k}$ 可表示为 $B_{k}\cap A_{k}$, 而 $B_{k}=(-\infty,b_{k})$, $A_{k}=(a_{k},\infty)$. 由于 $f$ 是可测函数, 每个 $f^{-1}(B_{k})$ 和 $f^{-1}(A_{k})$ 是可测集. 另一方面, 可测集是 $\sigma$ 代数, 因此 $f^{-1}(\mathcal{O})$ 是可测的, 这是由于
\begin{align*}
f^{-1}(\mathcal{O})=f^{-1}\left[\bigcup_{k=1}^{\infty}B_{k}\cap A_{k}\right]=\bigcup_{k=1}^{\infty}f^{-1}(B_{k})\cap f^{-1}(A_{k}).
\end{align*}
\end{proof}

以下命题告诉我们初等分析中最为熟悉的连续函数是可测的.

\begin{proposition}\label{proposition:ch03-连续可测}
在可测集定义域上连续的实值函数是可测的.
\end{proposition}

\begin{proof}
令函数 $f$ 在可测集 $E$ 上连续. 令 $\mathcal{O}$ 为开的. 由于 $f$ 是连续的, $f^{-1}(\mathcal{O})=E\cap\mathcal{U}$, 其中 $\mathcal{U}$ 是开的. 因此 $f^{-1}(\mathcal{O})$ 作为两个可测集的交是可测的. 从前一个命题得知 $f$ 是可测的.
\end{proof}

递增或递减的实值函数称为是单调的. 我们将下一个命题的证明留作练习（见习题 24）.

\begin{proposition}\label{proposition:ch03-单调可测}
定义在区间上的单调函数是可测的.
\end{proposition}

\begin{proposition}\label{proposition:ch03-a.e.可测}
令 $f$ 为 $E$ 上的扩充的实值函数.
\begin{enumerate}[(i)]
\item 若 $f$ 在 $E$ 上可测, 而在 $E$ 上 $f=g$ a.e., 则 $g$ 在 $E$ 上可测.
\item 对于 $E$ 的可测子集 $D$, $f$ 在 $E$ 上是可测的当且仅当 $f$ 在 $D$ 和 $E\sim D$ 上的限制是可测的.
\end{enumerate}
\end{proposition}

\begin{proof}
首先假设 $f$ 是可测的. 定义 $A=\{x\in E\mid f(x)\neq g(x)\}$. 观察到
\begin{align*}
\{x\in E\mid g(x)>c\}=\{x\in A\mid g(x)>c\}\cup[\{x\in E\mid f(x)>c\}\cap[E\sim A]].
\end{align*}
由于在 $E$ 上 $f=g$ a.e., $m(A)=0$. 因此 $\{x\in A\mid g(x)>c\}$ 是可测的, 这是由于它是零测度集的子集. 集合 $\{x\in E\mid f(x)>c\}$ 是可测的, 这是由于 $f$ 在 $E$ 上是可测的. 由于 $E$ 和 $A$ 都是可测的, 而可测集是一个代数, 集合 $\{x\in E\mid g(x)>c\}$ 是可测的. 为验证 (ii), 只要观察到对任何 $c$,
\begin{align*}
\{x\in E\mid f(x)>c\}=\{x\in D\mid f(x)>c\}\cup\{x\in E\sim D\mid f(x)>c\}
\end{align*}
且又一次利用可测集是一个代数这一事实.
\end{proof}

两个可测的扩充实值函数 $f$ 和 $g$ 的和 $f+g$ 在那些 $f$ 和 $g$ 取异号的无穷值的点不是恰当定义的. 假设 $f$ 和 $g$ 在 $E$ 上 a.e. 有限. 定义集合 $E_{0}$ 为 $E$ 中那些使得 $f$ 和 $g$ 都是有限点的集合. 若 $f+g$ 在 $E_{0}$ 上的限制是可测的, 则由前一个命题, $f+g$ 的任何到整个 $E$ 的延拓, 作为扩充的实值函数也是可测的. 这是我们认为 "两个 a.e. 有限可测函数的和是可测" 毫无歧义的道理. 类似的说明适用于乘积. 以下命题告诉我们实施在 a.e. 有限可测函数上的标准的代数运算仍然导出可测函数.

\begin{theorem}\label{theorem:ch03-代数可测}
令 $f$ 和 $g$ 为 $E$ 上的可测函数, 在 $E$ 上 a.e. 有限.
\begin{description}
\item[（线性）] 对任何 $\alpha$ 和 $\beta$,
\begin{align*}
\alpha f+\beta g\text{ 在 }E\text{ 上可测}.
\end{align*}
\item[（乘积）]
\begin{align*}
fg\text{ 在 }E\text{ 上可测}.
\end{align*}
\end{description}
\end{theorem}

\begin{proof}
根据以上的说明, 我们可以假设 $f$ 和 $g$ 在整个 $E$ 上有限. 若 $\alpha=0$, 则函数 $\alpha f$ 也是可测的. 若 $\alpha\neq 0$, 观察到对于数 $c$,
\begin{align*}
\text{若 }\alpha>0,\quad\{x\in E\mid\alpha f(x)>c\}=\{x\in E\mid f(x)>c/\alpha\},
\end{align*}
而
\begin{align*}
\text{若 }\alpha<0,\quad\{x\in E\mid\alpha f(x)>c\}=\{x\in E\mid f(x)<c/\alpha\}.
\end{align*}
因此 $f$ 的可测性蕴涵 $\alpha f$ 的可测性. 因此为证明线性仅须考虑 $\alpha=\beta=1$ 的情形.

对于 $x\in E$, 若 $f(x)+g(x)<c$, 则 $f(x)<c-g(x)$, 从而根据有理数集 $\mathbf{Q}$ 在 $\mathbf{R}$ 中的稠密性, 存在有理数 $q$ 使得
\begin{align*}
f(x)<q<c-g(x).
\end{align*}
因此,
\begin{align*}
\{x\in E\mid f(x)+g(x)<c\}=\bigcup_{q\in\mathbf{Q}}\{x\in E\mid g(x)<c-q\}\cap\{x\in E\mid f(x)<q\}.
\end{align*}
有理数是可数的. 因此 $\{x\in E\mid f(x)+g(x)<c\}$ 是可测的, 这是由于它是可测集的可数族的并. 因此 $f+g$ 是可测的.

为证明可测函数的乘积是可测的, 首先观察到
\begin{align*}
fg=\frac{1}{2}\left[(f+g)^{2}-f^{2}-g^{2}\right].
\end{align*}
因此, 由于我们已证明线性, 为证明两个可测函数的乘积是可测的, 仅须证明可测函数的平方是可测的. 对 $c\geqslant 0$,
\begin{align*}
\{x\in E\mid f^{2}(x)>c\}=\{x\in E\mid f(x)>\sqrt{c}\}\cup\{x\in E\mid f(x)<-\sqrt{c}\},
\end{align*}
而对于 $c<0$,
\begin{align*}
\{x\in E\mid f^{2}(x)>c\}=E.
\end{align*}
因此 $f^{2}$ 是可测的.
\end{proof}

初等分析中考虑过的许多函数性质, 包括连续性和可微性, 在函数的复合运算下依然成立. 然而, 可测函数的复合可以不是可测的.

\begin{example}\label{example:ch03-复合不可测}
存在两个可测实值函数, 每个定义在整个 $\mathbf{R}$ 上, 其复合不是可测的. 根据第 2 章的\cref{proposition:ch02-Cantor映射}, 存在定义在 $[0,1]$ 上的连续严格递增函数 $\psi$ 和 $[0,1]$ 的可测子集 $A$, 使得 $\psi(A)$ 是不可测的. 延拓 $\psi$ 为 $\mathbf{R}$ 映上 $\mathbf{R}$ 的连续严格递增函数. 函数 $\psi^{-1}$ 是连续的, 因此是可测的. 另一方面, $A$ 是可测集, 从而它的特征函数（其定义见 3.2 节）$\chi_{A}$ 是可测函数. 我们宣称复合函数 $f=\chi_{A}\circ\psi^{-1}$ 是不可测的. 事实上, 若 $I$ 是任何包含 $1$ 但不包含 $0$ 的开区间, 则它在 $f$ 下的原象是不可测集 $\psi(A)$.
\end{example}

尽管有该例子施加的阻碍, 但存在关于在复合下可测性成立的如下有用的命题（也见习题 11）.

\begin{proposition}\label{proposition:ch03-外复合可测}
令 $g$ 为定义在 $E$ 上的可测实值函数, 而 $f$ 为定义在整个 $\mathbf{R}$ 上的连续实值函数. 则复合 $f\circ g$ 是 $E$ 上的可测函数.
\end{proposition}

\begin{proof}
根据\cref{proposition:ch03-开集原象}, 实值函数是可测的当且仅当每个开集的原象是可测的. 令 $\mathcal{O}$ 为开的. 则
\begin{align*}
(f\circ g)^{-1}(\mathcal{O})=g^{-1}(f^{-1}(\mathcal{O})).
\end{align*}
由于 $f$ 是连续的且定义在开集上, 集合 $\mathcal{U}=f^{-1}(\mathcal{O})$ 是开的. 我们从函数 $g$ 的可测性推出 $g^{-1}(\mathcal{U})$ 是可测的. 因此原象 $(f\circ g)^{-1}(\mathcal{O})$ 是可测的, 因而复合函数 $f\circ g$ 是可测的.
\end{proof}

上面结果的一个直接的重要推论是: 若具有定义域 $E$ 的函数 $f$ 是可测的, 则 $|f|$ 是可测的, 且事实上对每个 $p>0$, 具有共同的定义域 $E$ 的 $|f|^{p}$ 是可测的.

对于具有共同定义域 $E$ 的函数的有限簇 $\{f_{k}\}_{k=1}^{n}$, 函数 $\max\{f_{1},\cdots,f_{n}\}$ 在 $E$ 上的定义为对 $x\in E$, $\max\{f_{1},\cdots,f_{n}\}(x)=\max\{f_{1}(x),\cdots,f_{n}(x)\}$. 以相同方式定义函数 $\min\{f_{1},\cdots,f_{n}\}$.

\begin{proposition}\label{proposition:ch03-maxmin可测}
对于具有共同定义域 $E$ 的可测函数的有限簇 $\{f_{k}\}_{k=1}^{n}$, 函数 $\max\{f_{1},\cdots,f_{n}\}$ 和 $\min\{f_{1},\cdots,f_{n}\}$ 也是可测的.
\end{proposition}

\begin{proof}
对任何 $c$, 我们有
\begin{align*}
\{x\in E\mid\max\{f_{1},\cdots,f_{n}\}(x)>c\}=\bigcup_{k=1}^{n}\{x\in E\mid f_{k}(x)>c\}.
\end{align*}
因此 max 函数是可测的. min 的情形类似.
\end{proof}


对于定义在 $E$ 上的函数 $f$, 我们有与之相联系的定义在 $E$ 上的函数 $|f|$、$f^{+}$ 和 $f^{-}$,
\begin{align*}
|f|(x)=\max\{f(x),-f(x)\},\quad f^{+}(x)=\max\{f(x),0\},\quad f^{-}(x)=\max\{-f(x),0\}.
\end{align*}
若 $f$ 在 $E$ 上可测, 则根据前一个命题, 函数 $|f|$、$f^{+}$ 和 $f^{-}$ 也可测. 当我们研究积分时这是重要的, 由于将 $f$ 表示为 $E$ 上的两个非负函数的差
\begin{align*}
f=f^{+}-f^{-}
\end{align*}
的表示方式在定义 Lebesgue 积分时起着重要作用.

\begin{exercise}\label{exercise:ch03-01}
假定 $f$ 和 $g$ 是 $[a,b]$ 上的连续函数. 证明: 若在 $[a,b]$ 上 $f=g$ a.e., 则事实上在 $[a,b]$ 上 $f=g$. 若 $[a,b]$ 换成一般的可测集 $E$, 类似的断言是否成立?
\end{exercise}

\begin{exercise}\label{exercise:ch03-02}
令 $D$ 和 $E$ 为可测集, 而 $f$ 是以 $D\cup E$ 为定义域的函数. 我们证明 $f$ 在 $D\cup E$ 上是可测的当且仅当它限制在 $D$ 和 $E$ 上是可测的. 若 "可测的" 换作 "连续的", 同样结论是否成立?
\end{exercise}

\begin{exercise}\label{exercise:ch03-03}
假定函数 $f$ 有可测定义域且除有限个点外是连续的. $f$ 是否必然可测?
\end{exercise}

\begin{exercise}\label{exercise:ch03-04}
假定 $f$ 是 $\mathbf{R}$ 上的实值函数, 使得对每个数 $c$, $f^{-1}(c)$ 是可测的. $f$ 必然是可测的吗?
\end{exercise}

\begin{exercise}\label{exercise:ch03-05}
假定函数 $f$ 定义在可测集 $E$ 上且有性质: 对每个有理数 $c$, $\{x\in E\mid f(x)>c\}$ 是可测集. $f$ 必然是可测函数吗?
\end{exercise}

\begin{exercise}\label{exercise:ch03-06}
令 $f$ 为具有可测定义域 $D$ 的函数. 证明 $f$ 是可测的当且仅当 $\mathbf{R}$ 上的定义如下的函数 $g$: 对 $x\in D$, $g(x)=f(x)$, 而对 $x\notin D$, $g(x)=0$ 是可测的.
\end{exercise}

\begin{exercise}\label{exercise:ch03-07}
令函数 $f$ 为定义在可测集 $E$ 上的函数. 证明 $f$ 是可测的当且仅当对每个 Borel 集 $A$, $f^{-1}(A)$ 是可测的.（提示: 具有性质 $f^{-1}(A)$ 是可测的集合 $A$ 的族是一个 $\sigma$ 代数.）
\end{exercise}

\begin{exercise}\label{exercise:ch03-08}
（Borel 可测性）函数 $f$ 称为是 Borel 可测的, 若它的定义域 $E$ 是 Borel 集, 且对每个 $c$, 集合 $\{x\in E\mid f(x)>c\}$ 是 Borel 集. 证明若我们用 "Borel 集" 代替 "Lebesgue 可测集", 命题 1 和定理 6 仍然成立. 证明:
\begin{enumerate}[(i)]
\item 每个 Borel 可测函数是 Lebesgue 可测的;
\item 若 $f$ 是 Borel 可测的且 $B$ 是 Borel 集, 则 $f^{-1}(B)$ 是一个 Borel 集;
\item 若 $f$ 和 $g$ 是 Borel 可测的, 则 $f\circ g$ 也是;
\item 若 $f$ 是 Borel 可测的且 $g$ 是 Lebesgue 可测的, 则 $f\circ g$ 是 Lebesgue 可测的.
\end{enumerate}
\end{exercise}

\begin{exercise}\label{exercise:ch03-09}
令 $\{f_{n}\}$ 为定义在可测集 $E$ 上的可测函数序列. 定义 $E_{0}$ 为 $E$ 中那些使得 $\{f_{n}(x)\}$ 收敛的点 $x$ 所组成的集合. 集合 $E_{0}$ 是可测的吗?
\end{exercise}

\begin{exercise}\label{exercise:ch03-10}
假定 $f$ 和 $g$ 是定义在整个 $\mathbf{R}$ 上的实值函数, $f$ 是可测的, 而 $g$ 是连续的. 复合 $f\circ g$ 必然是可测的吗?
\end{exercise}

\begin{exercise}\label{exercise:ch03-11}
令 $f$ 为可测函数而 $g$ 是从 $\mathbf{R}$ 映上 $\mathbf{R}$ 的具有 Lipschitz 逆的一对一的函数. 证明复合 $f\circ g$ 是可测的.（提示: 参考第 2 章的习题 37.）
\end{exercise}

\end{document}